%#!ratex
\documentclass[pdflatex,ja=standard,a5paper]{bxjsarticle}
\usepackage{amsmath,twemojis}
\newcommand*\cEmoji[1]{%
  \raisebox{-.12\zw}[.88\zw]{\twemoji[height=1\zw]{#1}}}
\title{ナベアツ総和方程式}
\author{某~ZR\cEmoji{upside_down_face}}
\begin{document}
\maketitle
\section{いきなり結論}
$d$~を正整数とする。%
$1$~から~$10^d$~の範囲にあるナベアツ数の総和~$S$~
は次の式で求められる。
\[
S = \begin{cases}
18 & \text{（}d = 1\text{~のとき）}\\
\dfrac{(10^d - 1)(81\cdot10^d - 56\cdot9^d)}{162} &
  \text{（それ以外）}
\end{cases}
\]
\section{まとめ}
もっと本質的\cEmoji{snowman}な文書をつくりましょう！%
\cEmoji{information_desk_person}
\end{document}
